\chapter{Software Engineering}
\label{ch:engenharia}

\section{Introduction}

This chapter details system architecture and design: functional and non-functional requirements,
diagrams, data model, and security architecture.

\section{System Overview}

The following figures summarise the end-to-end workflow and high-level architecture; the diagram
subsections below refine each element.

\introdiagram[End-to-end workflow (overview sequence)]{%
Figure~\ref{fig:intro-workflow} situates the analyst workflow: a sample enters through the web or
desktop interface, undergoes static analysis on the host, and may optionally be executed in an
isolated Hyper-V sandbox. Results converge in a structured report with decompiled code and an
explainable risk score.}{fig-4-5-sequence-overview.png}{End-to-end workflow---sample submission, static analysis, optional sandbox execution, and unified reporting (overview sequence)}{fig:intro-workflow}

\introdiagram[High-level system pipeline]{%
Figure~\ref{fig:intro-pipeline} shows how the presentation tier, API, analysis engine, and sandbox
layers relate at architectural level.}{fig-4-4-architecture.png}{High-level system pipeline---presentation, API, analysis engine, and sandbox isolation}{fig:intro-pipeline}

\section{Requirements Analysis}

This section formalises what \rat\ must deliver and the quality attributes it must uphold.
Functional capabilities are tabulated first, non-functional constraints second, and
traceability to implementation components and automated tests last. Elicitation drew on the
analyst workflow from Chapter~\ref{ch:intro}, supervisor feedback, and positioning against
existing sandboxes (Chapter~\ref{ch:estado-arte}).

\subsection{Functional Requirements}

Functional requirements were elicited from the intended analyst workflow, the need for isolated
dynamic execution, and the academic goal of explainable triage rather than opaque classification.
Stakeholder interviews with the project supervisor, together with comparison against commercial
sandboxes and open-source frameworks (Chapter~\ref{ch:estado-arte}), informed the scope of upload
validation, analysis modes, and reporting features. Table~\ref{tab:fr} presents the functional
requirements identified for \rat.

\small
\setlength{\LTleft}{0pt}
\setlength{\LTright}{0pt}
\begin{xltabular}{\linewidth}{@{}>{\raggedright\arraybackslash}p{1.5cm}Y@{}}
\caption{Functional requirements of \rat} \label{tab:fr} \\
\toprule
\textbf{ID} & \textbf{Description} \\
\midrule
\endfirsthead
\toprule
\textbf{ID} & \textbf{Description} \\
\midrule
\endhead
\fr{01} &
The system shall accept sample uploads in \texttt{.exe}, \texttt{.dll}, and \texttt{.cs} formats
via the web interface, desktop application, or REST API. The system shall validate file extension,
enforce a configurable size limit (default 100~MB), and sanitise file names. For \texttt{.cs}
submissions on the web/API, the system shall apply string-based heuristics only; full PE analysis
shall require a compiled binary. \\
\fr{02} &
The system shall execute a seven-phase static analysis pipeline comprising PE metadata extraction,
managed decompilation when applicable, import/string/evasion heuristics, deobfuscation, YARA
scanning, risk scoring, native decompilation when required, and structured report generation. \\
\fr{03} &
The system shall execute samples in an isolated Hyper-V virtual machine without Internet
connectivity and return a behavioural report. The user shall select static, dynamic, or combined
analysis modes; the web interface shall use Path~A and the WPF application shall use Path~B. \\
\fr{04} &
The system shall create, query, and list analysis jobs identified by UUID. The system shall persist
job metadata in SQLite and store per-job artefacts. \\
\fr{05} &
The system shall present analysis results in a web interface with report, pseudo-C, and IL panels.
The user shall navigate cross-references between panels. The system shall provide a permalink with
live status updates during sandbox runs. \\
\fr{06} &
The system shall compute a risk score from 0 to 100 by aggregating YARA matches, C2 indicators,
evasion techniques, packers, obfuscation, persistence signals, and entropy. The system shall
assign one of five severity levels and apply strong-evidence gates for the highest tiers. \\
\fr{07} &
The system shall stream incremental NDJSON log events during static analysis so that the user
receives real-time progress feedback. \\
\fr{08} &
The system shall orchestrate dynamic analysis through two paths: (A) a backend-driven HTTP agent
inside the guest, and (B) PowerShell Direct with post-transfer integrity verification from the WPF
client. Path~B shall publish sandbox results to the same job as a prior static run when applicable. \\
\fr{09} &
The system shall offer an optional browser-side assistant that explains visible decompiled
pseudo-C excerpts. The user's Gemini API key shall remain in local browser storage. The system
shall provide a demonstration mode that does not contact external services. \\
\bottomrule
\end{xltabular}
\normalsize

\subsection{Non-Functional Requirements}

Non-functional requirements express quality attributes essential to malware analysis in a
laboratory environment: confidentiality of samples, isolation of execution, usability for students,
and maintainability for long-term academic use. Table~\ref{tab:nfr} lists the non-functional
requirements alongside their rationale.

\small
\setlength{\LTleft}{0pt}
\setlength{\LTright}{0pt}
\begin{xltabular}{\linewidth}{@{}>{\raggedright\arraybackslash}p{1.5cm}Y@{}}
\caption{Non-functional requirements of \rat} \label{tab:nfr} \\
\toprule
\textbf{ID} & \textbf{Description} \\
\midrule
\endfirsthead
\toprule
\textbf{ID} & \textbf{Description} \\
\midrule
\endhead
\nfr{01} &
The system shall bind the API to localhost in production deployments. The system shall require
authentication tokens before serving analysis requests. Secrets shall not be stored in version
control. \\
\nfr{02} &
The system shall connect the sandbox virtual machine to an internal switch without Internet
routing. The system shall restore a clean snapshot after each dynamic run. Path~B shall verify
network isolation before sample execution. \\
\nfr{03} &
The system shall support drag-and-drop upload, real-time progress display, and a responsive layout.
A user with basic security knowledge shall be able to interpret results without formal training. \\
\nfr{04} &
The system shall support at least two concurrent analysis workers by default. The administrator
shall configure execution and VM-operation timeouts per sandbox path. \\
\nfr{05} &
The system shall be verifiable through automated continuous integration (pytest, Vitest, dotnet
build). Dependencies shall be locked and technical documentation versioned. \\
\nfr{06} &
The system shall operate without an external Redis instance by falling back to local worker threads.
An optional Redis queue shall be available for distributed deployments. \\
\bottomrule
\end{xltabular}
\normalsize

\subsection{Requirements Traceability}

Each functional requirement maps to a primary component (backend API, static pipeline, sandbox,
job store, web UI, scoring module, streaming API, VM orchestration, Gemini assistant) and is
verified by automated tests in continuous integration. The full traceability matrix, listing
modules and test suites per requirement, is omitted here to avoid duplicating repository structure;
coverage is demonstrated by the 220 automated tests summarised in Section~\ref{sec:testes}.

\FloatBarrier
\section{Diagrams}

PlantUML diagrams (synchronised via \filepath{sync\_diagrams\_to\_report.py}) document functional
flows, architecture, security boundaries, and sandbox orchestration. Each subsection below
introduces one diagram and explains its role in the design; the figures are referenced throughout
Chapters~\ref{ch:engenharia} and~\ref{ch:implementacao}.

\diagramsubsection{Use Case Diagram}{%
Figure~\ref{fig:usecase} identifies the primary actors---analyst, backend services, sandbox guest,
and optional LLM assistant---and the use cases they trigger, from sample upload through static
analysis, dynamic execution, and report consultation.}{fig-4-1-usecase.png}{Use case diagram - \rat}{fig:usecase}

\diagramsubsection{Activity Diagram - Static Analysis}{%
Figure~\ref{fig:activity-static} details the seven-phase static pipeline executed on the host:
metadata extraction, managed and native decompilation branches, heuristic scanning, YARA matching,
risk scoring, and report assembly.}{fig-4-2-activity-static.png}{Activity diagram - static analysis pipeline}{fig:activity-static}

\diagramsubsection{Activity Diagram - Dynamic Analysis}{%
Figure~\ref{fig:activity-dynamic} shows the dynamic path from job submission through snapshot
restore, guest execution, telemetry collection, and report transfer back to the host.}{fig-4-3-activity-dynamic.png}{Activity diagram - dynamic analysis pipeline}{fig:activity-dynamic}

\diagramsubsection{System Architecture}{%
Figure~\ref{fig:architecture} presents the conceptual layering of \rat: presentation tier (web and
WPF), FastAPI backend, analysis modules, job store, and Hyper-V sandbox.}{fig-4-4-architecture.png}{Conceptual architecture of \rat}{fig:architecture}

\introdiagram{%
Figure~\ref{fig:sequence} extends this view with an integrated sequence diagram covering static,
dynamic, and WPF-initiated flows on a single timeline.}{fig-4-5-sequence-overview.png}{Integrated sequence diagram - static, dynamic, and WPF flows}{fig:sequence}

\diagramsubsection{Trust Zones and Authentication}{%
Figure~\ref{fig:trust-zones} maps the three trust zones---host, isolated sandbox network, and
untrusted guest---and the data flows crossing each boundary.}{fig-4-6-security-trust-zones.png}{Trust zones - host / sandbox / guest isolation}{fig:trust-zones}

\introdiagram{%
Figure~\ref{fig:auth} complements the trust-zone view with the authentication tokens and
credentials required by each component.}{fig-4-10-security-auth.png}{Authentication map - tokens and credentials}{fig:auth}

\diagramsubsection{Sandbox Paths (A and B)}{%
Figure~\ref{fig:sandbox-paths} compares the two dynamic orchestration paths at a glance.}{fig-4-7-sandbox-paths.png}{Comparison of paths A (HTTP vm-agent) and B (Hyper-V PowerShell)}{fig:sandbox-paths}

\introdiagram{%
Figure~\ref{fig:seq-path-a} sequences the HTTP vm-agent interactions on Path~A: snapshot restore,
sample upload, execution, and report retrieval.}{fig-4-8-sequence-path-a.png}{Sequence - Path A (HTTP vm-agent)}{fig:seq-path-a}

\introdiagram{%
Figure~\ref{fig:seq-path-b} details the PowerShell Direct workflow on Path~B, including SHA-256
integrity verification of the report copied from the guest.}{fig-4-9-sequence-path-b.png}{Sequence - Path B (PsDirect + SHA256 verification)}{fig:seq-path-b}

\FloatBarrier
\section{Security Architecture}
\label{sec:seguranca}

Three trust zones apply: \textbf{host} (trusted analysis services), \textbf{sandbox network}
(isolated internal switch), and \textbf{guest} (untrusted sample execution). The design follows a
\emph{fail-closed} policy: production services refuse to start without configured authentication
tokens; malware executes only inside the VM or a non-executing stub driver; secrets remain outside
version control; the API listens on localhost; and guest accounts operate with least privilege while
host administrators manage Hyper-V lifecycle. Defence in depth combines token authentication, upload
rate limiting, name sanitisation, size caps, and mandatory snapshot restore between runs.

Authentication is layered by component: the backend and WPF client share an API token; the in-guest
agent validates a separate agent token; Hyper-V automation uses PowerShell Direct with dedicated
guest credentials. The optional Gemini integration never receives backend credentials---the user's
API key stays in browser storage---reducing the attack surface for third-party LLM providers.
Primary residual risks include credential exposure in PowerShell logs (mitigated by environment
variables and secure dialogs), optional client-side code exfiltration to Gemini (user-controlled),
and dependency tampering (mitigated by hash-locked installs and optional binary verification on
WPF startup). Figures~\ref{fig:trust-zones} and~\ref{fig:auth} illustrate these zones and token
flows; detailed endpoint and variable listings appear in the appendix for operators.

\section{Sandbox Path Comparison}
\label{sec:caminhos-ab}

Paths A and B address the same isolation goal through different orchestration trade-offs
(Figure~\ref{fig:sandbox-paths}). \textbf{Path~A} integrates dynamic analysis into the Python
backend. After snapshot restore, the host talks to an HTTP agent inside the guest for upload,
execution, and report retrieval. This path suits the web workflow but collects lighter telemetry.
\textbf{Path~B} is driven by the WPF desktop and PowerShell. Samples reach the guest via
PowerShell Direct or VM file copy; a guest-side collector gathers richer diffs (files, registry,
network, Sysmon when present); and the host verifies report integrity with SHA-256 before accepting
results. Recent sandbox systematisation work \cite{alrawi2024} stresses that monitoring depth and
configuration strongly affect observable behaviour---a finding that motivated keeping Path~B as a
deeper, independently evolvable instrumentation path. Both paths share Hyper-V, the internal switch,
and clean-snapshot restore.

\section{Data Model}

Persistence relies on \textbf{SQLite} with one record per analysis job (UUID identifier, analysis
type, status, file metadata, timestamps). Per-job artefacts include static and dynamic reports,
deobfuscation excerpts, and an optional sample copy. Jobs progress through four states---\emph{queued},
\emph{running}, \emph{completed}, and \emph{failed}---enabling the web interface to poll status
during long sandbox runs.

Table~\ref{tab:scoring}
summarises the maximum caps per factor in the risk scorer (each factor contributes up to
its respective cap; the final score is limited to 100).

\begin{table}[htbp]
    \centering
    \caption{Risk score factors (0--100)}
    \label{tab:scoring}
    \begin{tabularx}{\linewidth}{@{}Yr@{\hspace{0.6cm}}Yr@{}}
        \toprule
        \textbf{Factor} & \textbf{Cap} & \textbf{Factor} & \textbf{Cap} \\
        \midrule
        C2 strings & 20 & Persistence & 10 \\
        YARA matches & 25 & Packer & 10 \\
        Suspicious functions & 20 & Obfuscation & 8 \\
        Suspicious imports & 12 & High entropy & 5 \\
        Stealer indicators & 15 & & \\
        Evasion techniques & 12 & & \\
        \bottomrule
    \end{tabularx}
\end{table}

\textbf{Levels}: CRITICAL if score $\geq 72$ and
\texttt{strong\_evidence} (YARA match, or strong C2 combined with suspicious functions/imports);
HIGH if score $\geq 48$ with at least one strong signal; MEDIUM if score $\geq 35$; LOW if
score $\geq 15$; otherwise VERY LOW. (The production API returns Portuguese labels:
\texttt{CRÍTICO}, \texttt{ALTO}, \texttt{MÉDIO}, \texttt{BAIXO}, \texttt{MUITO BAIXO}.)

\section{Conclusions}

Modelling separates presentation, API, analysis, and sandbox concerns. Requirements remain
traceable through automated CI coverage rather than exhaustive in-report module listings.
